\chapter{总结与展望}

\section{工作总结}
本文围绕变电站运维场景下移动机器人自主导航的实际需求，针对复杂环境感知、定位建图、动态障碍物处理以及路径规划等关键问题，开展了面向变电站运维的搬运机器人导航系统设计与实现研究。全文主要工作总结如下：

\begin{enumerate}[label=(\arabic*)]
\item  面向变电站高安全性、高可靠性和复杂空间约束等应用特点，本文分析了运维机器人的功能需求，设计了机器人硬件平台和导航系统总体框架。系统以激光雷达、双目相机等传感器为环境信息输入，结合中央控制模块和底层运动控制模块，为后续环境感知、地图构建、路径规划和运动控制提供了整体实现基础。

\item  针对导航系统中多传感器数据统一表达的问题，本文梳理了世界坐标系、机器人基坐标系、雷达坐标系和相机坐标系之间的转换关系，并介绍了K-Means和DBSCAN等点云聚类基础方法。上述理论为点云处理、障碍物提取和地图构建提供了基础支撑。

\item  针对复杂动态环境中的障碍物感知问题，本文设计了基于深度图像和点云数据融合的轻量级障碍物检测与跟踪方法。该方法结合U深度图检测器和DBSCAN检测器的优势，通过包围盒匹配和融合降低单一传感器噪声带来的误检影响；同时引入基于特征的数据关联、恒加速度模型和点云投票策略，对动态障碍物进行识别与状态估计。在地图构建方面，本文采用局部占用栅格更新，为规划器提供可用于避障决策的环境约束。

\item  针对变电站设备密集、障碍形态复杂且可能存在动态干扰的导航场景，本文设计了全局参考引导与局部滚动优化相结合的路径规划方法。全局层采用目标偏向RRT生成二维拓扑可行参考路径，并通过直连检查、线段碰撞检测、超时处理和路径平滑复检提高工程可用性；局部层构建学习增强MPC方法，在承接障碍物边界和恒加速度状态估计结果的基础上，引入点级障碍表示、动态预测点流和学习型距离特征，降低非凸障碍包络造成的可行空间损失。本文同时完成了规划输出到全向底盘速度接口的设计，使滚动优化得到的首个控制量能够直接转化为机器人可执行指令。
\end{enumerate}

\section{工作展望}
本文围绕变电站运维机器人导航系统的关键模块进行了设计与研究，但距离工程化长期稳定运行仍有进一步完善空间。后续工作可从以下几个方面展开：

\begin{enumerate}[label=(\arabic*)]
\item  进一步完善真实变电站或高保真仿真场景下的系统实验。后续可围绕定位误差、建图精度、障碍物检测准确率、规划耗时、路径平滑性和任务完成率等指标建立更完整的评价体系，并开展多场景、多工况对比实验。

\item  提升复杂电力场景下的多传感器融合鲁棒性。变电站中金属结构密集、电磁环境复杂，单一传感器容易受到遮挡、反射或噪声影响。后续可进一步融合激光雷达、视觉、IMU等多源信息，提高机器人在弱纹理、强反光、遮挡和动态干扰场景下的定位与感知稳定性。

\item  优化动态障碍物预测与局部规划策略。当前方法主要基于显式运动模型和几何约束进行动态避障，后续可结合学习型运动预测、风险场建模或神经规划方法，提高机器人在非凸障碍物、突发动态目标和狭窄通道中的通行能力。

\item  推进导航系统与具体运维任务的深度结合。后续可将导航系统与设备巡检、仪表识别、目标抓取、辅助搬运等任务模块进一步集成，形成从任务分配、路径规划、环境感知到动作执行的闭环作业系统，从而提升变电站机器人系统的实用价值。
\end{enumerate}
